\documentclass[11pt]{amsart}
\usepackage{amsmath,amssymb,amsthm}
\usepackage[margin=1.15in]{geometry}
\usepackage[hidelinks]{hyperref}
\newtheorem{theorem}{Theorem}[section]
\newtheorem{corollary}[theorem]{Corollary}
\newtheorem{lemma}[theorem]{Lemma}
\theoremstyle{remark}
\newtheorem{remark}[theorem]{Remark}
\newtheorem{example}[theorem]{Example}
\newcommand{\doi}[1]{\href{https://doi.org/#1}{doi:\nolinkurl{#1}}}
\newcommand{\arxiv}[1]{\href{https://arxiv.org/abs/#1}{arXiv:\nolinkurl{#1}}}
\newcommand{\rk}{\operatorname{rk}}
\setlength{\emergencystretch}{2.5em}

\begin{document}

\title[Coalescence of grand couplings]{Almost every finite Markov chain has only coalescing grand couplings}
\author{Yitzchak Shmalo}
\address{Einstein Institute of Mathematics, The Hebrew University of Jerusalem, Jerusalem, Israel}
\email{yspennstate@gmail.com}
\subjclass[2020]{60J10, 60J22, 15B51, 20M20}
\keywords{Markov chain, grand coupling, coalescence, coupling from the past, doubly stochastic matrix,
transformation semigroup}

\begin{abstract}
A grand coupling of a finite Markov chain runs copies of the chain from every state by applying independent random
maps with a common law; its coalescence number is the eventual number of distinct trajectories. Grimmett and Holmes
asked whether, for a random transition matrix with independent uniform weights, every grand coupling has
coalescence number one almost surely. We show that it does. If some grand coupling of $P$ has coalescence number
$k\ge2$, then there are $k$ states $W$ and a partition $\{Z_w:w\in W\}$ of the state space with $w\in Z_w$ such that the
matrix $\bigl(P(w,Z_{w'})\bigr)_{w,w'\in W}$ is doubly stochastic; for $k$ equal to the number of states this is the
Birkhoff--von Neumann condition. Each such condition is a non-trivial affine equation in $P$, and there are finitely
many, so the matrices with a non-coalescing grand coupling form a Lebesgue-null set.
\end{abstract}

\maketitle

\section{Introduction}\label{sec:intro}

Let $S=\{1,\dots,n\}$ with $n\ge2$, and let $P=(p_{i,j})$ be an irreducible transition matrix on $S$. Write $F_S$ for
the set of maps $f:S\to S$. Following Grimmett and Holmes~\cite{GH21,GH26}, a \emph{grand coupling} of $P$ is a
probability measure $\mu$ on $F_S$ such that
\begin{equation}\label{eq:consistent}
\mu\{f\in F_S: f(i)=j\}=p_{i,j}\qquad\text{for all }i,j\in S.
\end{equation}
Let $L_P$ be the set of grand couplings of $P$. Given $\mu\in L_P$, let $F_1,F_2,\dots$ be independent random maps with
law $\mu$ and put $X^i_t=F_t\circ\dots\circ F_1(i)$. Each $X^i$ is a Markov chain with transition matrix $P$ started at
$i$, and two of these chains move together once they meet. The number of coalescence classes, that is, of
trajectories that remain distinct for all time, is almost surely constant~\cite[Lemma~2]{GH21}; it is the
\emph{coalescence number} $k(\mu)$, and
\[
K(P)=\{k(\mu):\mu\in L_P\}.
\]
The coalescence of all trajectories is what the coupling-from-the-past algorithm of Propp and Wilson~\cite{PW96}
requires, and Grimmett and Holmes showed that the forward and backward processes have the same set
$K(P)$~\cite{GH21}. By~\cite[Theorem~3.13]{GH26}, $1\in K(P)$ if and only if $P$ is aperiodic, the independence coupling
attaining $k=1$. At the other end, $n\in K(P)$ if and only if $P$ is doubly stochastic~\cite[Theorem~4]{GH21}: a
coupling with $k(\mu)=n$ is supported on permutations, and conversely a doubly stochastic matrix is a mixture of
permutation matrices by the theorem of Birkhoff and von Neumann~\cite{Bir46} (see~\cite[Theorem~2.6]{GH26}), so the
coupling that picks a random permutation with these weights has $k=n$.

To ask what happens for a typical matrix, Grimmett and Holmes~\cite[Remark~2.4]{GH26} let $Q$ be the law of the matrix
with entries $p_{i,j}=q_{i,j}/\sum_{l}q_{i,l}$, where the $q_{i,j}$ are independent and uniform on $(0,1)$. A $Q$-random
matrix has all entries in $(0,1)$, so it is irreducible and aperiodic, and it has many grand couplings. From a bound
on the largest coalescence number in terms of the invariant distribution they deduced that $K(P)=\{1\}$ with positive
$Q$-probability~\cite[Theorem~3.5, Remark~3.7]{GH26}, and they showed that for $Q$-a.e.\ $P$ every block measure is
trivial~\cite[Lemma~5.1]{GH26}. They asked~\cite[Open problem~3.8]{GH26}: is $K(P)=\{1\}$ for $Q$-a.e.\ $P$? We answer
this question.

\begin{theorem}\label{thm:main}
Let $n\ge2$. The set of irreducible transition matrices $P$ on $S$ with $K(P)\neq\{1\}$ is contained in a finite union
of affine hyperplanes, none of which contains the set of stochastic matrices. In particular it has Lebesgue measure
zero in the set of stochastic matrices, and $K(P)=\{1\}$ for $Q$-a.e.\ $P$; the same holds for any law of $P$ that is
absolutely continuous with respect to Lebesgue measure.
\end{theorem}

Theorem~\ref{thm:main} follows from a necessary condition for $k\in K(P)$, which may be of independent interest.

\begin{theorem}\label{thm:structure}
Let $P$ be an irreducible transition matrix on $S$ and let $\mu\in L_P$ have $k=k(\mu)\ge2$. Then there exist a set
$W\subseteq S$ with $|W|=k$ and a partition $\{Z_w:w\in W\}$ of $S$ with $w\in Z_w$ for each $w$, such that
\begin{equation}\label{eq:columns}
\sum_{x\in W}P(x,Z_w)=1\qquad\text{for every }w\in W,
\end{equation}
where $P(x,Z)=\sum_{z\in Z}p_{x,z}$. Equivalently, the $k\times k$ matrix $B=\bigl(P(w,Z_{w'})\bigr)_{w,w'\in W}$ is
doubly stochastic. Moreover the pair $(W,\{Z_w\})$ can be chosen to depend only on the support of $\mu$, and
\eqref{eq:columns} holds with $P(x,Z_w)$ replaced by $\mathbf{1}[g(x)\in Z_w]$ for every map $g$ in that support.
\end{theorem}

The rows of $B$ sum to one because the sets $Z_{w'}$ partition $S$, so \eqref{eq:columns} is exactly the statement
that the columns of $B$ sum to one. When $k=n$ we have $W=S$ and $Z_w=\{w\}$, and \eqref{eq:columns} says that $P$ is
doubly stochastic, which recovers the easy half of the Birkhoff--von Neumann characterization of $n\in K(P)$. In
general $W$ is the image of an idempotent map of least rank in the semigroup generated by the support of $\mu$, and
$Z_w$ is the fibre of that idempotent over $w$; the condition says that each step of the chain, followed by this
idempotent, permutes $W$. Irreducibility is assumed to match the setting of~\cite{GH26}; the proofs below do not use
it, so both theorems hold for every transition matrix on $S$.

The question, the model, the coalescence number and the law $Q$ are due to Grimmett and Holmes, who also showed that
$k(\mu)$ is the least rank in the semigroup generated by the support of $\mu$~\cite[Theorem~3(a)]{GH21}. The
permutation property behind Theorem~\ref{thm:structure} is not new either: for a word of least rank, Yano~\cite[\S4.2,
(4.6)--(4.9)]{Yan13} shows that every map in the support, followed by that word, permutes its image, and he uses the
resulting permutation process to study random walks driven by road colorings. The images of least size and the fact
that every word maps one of them onto another go back to Friedman's work on the road coloring problem~\cite{Fri90}
(see~\cite[Lemma~2.3]{Yan13}), and Budzban and Mukherjea~\cite{BM00} study that problem through the semigroup
generated by the colors. What is new here is the step from this pointwise statement to a linear condition on $P$ itself,
Theorem~\ref{thm:structure}. Its consequence is Theorem~\ref{thm:main}: the obstruction to $k\ge2$ is always one of
finitely many affine equations in the entries of $P$, whatever the support of the coupling, and this answers Open
problem~3.8 of~\cite{GH26}.

Section~\ref{sec:semigroup} collects the facts about transformation semigroups, Section~\ref{sec:proofs} proves the two
theorems, and Section~\ref{sec:checks} describes a computer check for $n\le5$.

\section{Least rank in a semigroup of maps}\label{sec:semigroup}

For $f\in F_S$ write $\rk f=|f(S)|$. For a finite set $\Sigma\subseteq F_S$ let $M(\Sigma)$ be the set of all
compositions $g_t\circ\dots\circ g_1$ with $t\ge0$ and $g_1,\dots,g_t\in\Sigma$ (the empty composition is the identity),
and let $r(\Sigma)=\min\{\rk f:f\in M(\Sigma)\}$.

\begin{lemma}[{\cite[Theorem~3(a)]{GH21}}]\label{lem:rank}
Let $\mu\in L_P$ and $\Sigma=\operatorname{supp}\mu$. Then $k(\mu)=r(\Sigma)$.
\end{lemma}

\begin{proof}
Since $F_S$ is finite, so is $\Sigma$, and every element of $\Sigma$ has positive $\mu$-mass. Let
$A_t=F_t\circ\dots\circ F_1(S)$, the set of positions occupied at time $t$; the number of distinct trajectories at time
$t$ is $|A_t|$, which is non-increasing in $t$, and $k(\mu)=\lim_t|A_t|$. Each $F_t\circ\dots\circ F_1$ lies in
$M(\Sigma)$, so $|A_t|\ge r(\Sigma)$. Fix $h=g_L\circ\dots\circ g_1\in M(\Sigma)$ with $\rk h=r(\Sigma)$. Almost surely
there is a time $t$ with $(F_{t+1},\dots,F_{t+L})=(g_1,\dots,g_L)$, and then
$A_{t+L}=h(A_t)\subseteq h(S)$, so $|A_{t+L}|\le r(\Sigma)$. Hence $|A_t|=r(\Sigma)$ for all large $t$.
\end{proof}

\begin{lemma}\label{lem:idempotent}
Let $\Sigma\subseteq F_S$ be finite and $r=r(\Sigma)$. There is $e\in M(\Sigma)$ with $\rk e=r$ and $e\circ e=e$. For such
$e$, let $W=e(S)$. Then $e$ fixes every point of $W$, and for every $g\in M(\Sigma)$ the map $x\mapsto e(g(x))$ is a
bijection of $W$ onto itself.
\end{lemma}

\begin{proof}
Take $f\in M(\Sigma)$ with $\rk f=r$. The powers of $f$ lie in the finite set $M(\Sigma)$, so $f^a=f^{a+b}$ for some
$a,b\ge1$, and then $e=f^{ab}$ satisfies $e\circ e=e$. Since $e(S)\subseteq f(S)$ and $\rk e\ge r=\rk f$, we have
$e(S)=f(S)$ and $\rk e=r$. Every $w\in W$ has the form $w=e(y)$, so $e(w)=e(e(y))=w$. Let $g\in M(\Sigma)$. Then
$e\circ g\circ e\in M(\Sigma)$, so its image $e(g(e(S)))=e(g(W))$ has at least $r=|W|$ elements. As $e(g(W))\subseteq W$,
the map $x\mapsto e(g(x))$ sends $W$ onto $W$, and it is a bijection because $W$ is finite.
\end{proof}

Lemma~\ref{lem:idempotent} is the familiar structure of the minimal ideal of a finite semigroup, specialised to
semigroups of maps; see~\cite[Chapter~1]{How95} for idempotent powers. Its last statement, for maps $g$ in the support
and with a word of least rank in place of $e$, is~\cite[(4.7)--(4.9)]{Yan13}; an idempotent is convenient here because
its fibres contain the points of its image. We also record which maps can carry a
coupling with $k\ge2$. For a non-empty family $\mathcal H$ of $k$-element subsets of $S$, let $\operatorname{End}(\mathcal H)$
be the set of $f\in F_S$ that map every member of $\mathcal H$ bijectively onto a member of $\mathcal H$, and let $M_f$
be the $0$--$1$ matrix of $f$, with $(M_f)_{i,j}=\mathbf 1[f(i)=j]$.

\begin{lemma}\label{lem:support}
Let $\mu\in L_P$ with $k=k(\mu)\ge2$. There is a non-empty family $\mathcal H$ of $k$-element subsets of $S$ with
$\operatorname{supp}\mu\subseteq\operatorname{End}(\mathcal H)$, and then $P=\sum_f\mu(f)M_f$ lies in the convex hull of
$\{M_f:f\in\operatorname{End}(\mathcal H)\}$.
\end{lemma}

\begin{proof}
Let $\Sigma=\operatorname{supp}\mu$, take $g_0\in M(\Sigma)$ with $\rk g_0=k$ (Lemma~\ref{lem:rank}) and put
$\mathcal H=\{h(g_0(S)):h\in M(\Sigma)\}$. Each member is the image of $h\circ g_0\in M(\Sigma)$, so it has at least $k$
elements, and at most $k$ since it is the image of a $k$-element set. For $f\in\Sigma$ and $A=h(g_0(S))\in\mathcal H$
we have $f(A)=(f\circ h)(g_0(S))\in\mathcal H$, so $|f(A)|=k=|A|$ and $f$ maps $A$ bijectively onto a member of
$\mathcal H$. The last statement is~\eqref{eq:consistent}.
\end{proof}

\section{Proofs of the theorems}\label{sec:proofs}

\begin{proof}[Proof of Theorem~\ref{thm:structure}]
Let $\Sigma=\operatorname{supp}\mu$. By Lemma~\ref{lem:rank}, $r(\Sigma)=k\ge2$. Take $e$ and $W=e(S)$ as in
Lemma~\ref{lem:idempotent}, so $|W|=k$, and set $Z_w=e^{-1}(w)$ for $w\in W$. These sets partition $S$, and $w\in Z_w$
because $e(w)=w$. Let $g\in\Sigma$ and $w\in W$. By Lemma~\ref{lem:idempotent} exactly one $x\in W$ has $e(g(x))=w$,
that is, $g(x)\in Z_w$. Thus
\begin{equation}\label{eq:pointwise}
\sum_{x\in W}\mathbf{1}[g(x)\in Z_w]=1\qquad\text{for every }g\in\Sigma\text{ and }w\in W.
\end{equation}
Integrating~\eqref{eq:pointwise} against $\mu$ and using~\eqref{eq:consistent} gives
$\sum_{x\in W}\mu\{g:g(x)\in Z_w\}=\sum_{x\in W}P(x,Z_w)=1$, which is~\eqref{eq:columns}. The pair $(W,\{Z_w\})$ was
built from $\Sigma$ alone.
\end{proof}

\begin{proof}[Proof of Theorem~\ref{thm:main}]
Suppose $K(P)\ne\{1\}$, and take $\mu\in L_P$ with $k(\mu)\ge2$. Theorem~\ref{thm:structure} gives a set $W$ with
$|W|\ge2$, a point $w\in W$ and a set $Z=Z_w$ such that
\begin{equation}\label{eq:hyperplane}
\sum_{x\in W}\sum_{z\in Z}p_{x,z}=1 .
\end{equation}
Here $w\in Z$, and $Z\ne S$ because any $w'\in W\setminus\{w\}$ lies in $Z_{w'}$, which is disjoint from $Z$. For fixed $n$
there are finitely many triples $(W,w,Z)$ with $2\le|W|$, $w\in W$ and $w\in Z\subseteq S\setminus(W\setminus\{w\})$, so
the set of such $P$ lies in a finite union of affine hyperplanes~\eqref{eq:hyperplane}.

It remains to show that each of these hyperplanes meets the set of stochastic matrices in a set of measure zero, for
Lebesgue measure on that set (a product of $n$ simplices) and for $Q$. Fix $(W,w,Z)$ and $x_1\in W$. Moving mass in row
$x_1$ from a state of $S\setminus Z$ to a state of $Z$, with the other rows fixed, changes the left side
of~\eqref{eq:hyperplane} and keeps $P$ stochastic, so the hyperplane does not contain the set of stochastic matrices, and
its intersection with the product of simplices, a convex set of full dimension, has measure zero. For $Q$, condition on
the rows $(q_{x,\cdot})_{x\ne x_1}$. Writing $c=1-\sum_{x\in W\setminus\{x_1\}}P(x,Z)$, which depends only on these rows,
equation~\eqref{eq:hyperplane} becomes
\[
(1-c)\sum_{z\in Z}q_{x_1,z}-c\sum_{z\notin Z}q_{x_1,z}=0 .
\]
The coefficients $1-c$ and $-c$ are not both zero, so this is a non-trivial linear equation in the vector
$(q_{x_1,z})_{z\in S}$, which is uniform on $(0,1)^n$ and independent of the other rows. Its conditional probability is
therefore zero, and by Fubini's theorem~\eqref{eq:hyperplane} has $Q$-probability zero. A finite union of null sets is
null. The same argument applies to any law of $P$ that is absolutely continuous with respect to Lebesgue measure on the
product of simplices.
\end{proof}

\begin{remark}\label{rem:necessary}
For $k=n$ the condition of Theorem~\ref{thm:structure} is also sufficient, by the Birkhoff--von Neumann theorem, as
recalled in the introduction. For $2\le k<n$ it is not sufficient in general. Take $n=3$ and $k=2$. For each choice of
$W$ and of the partition, the two column conditions in~\eqref{eq:columns} are equivalent, since the entries of $B$ sum
to $k$; so the condition defines finitely many pieces of hyperplanes, each meeting the interior of the set of stochastic
matrices (for instance $W=\{1,2\}$, $Z_1=\{1,3\}$, $Z_2=\{2\}$ gives $p_{1,2}+p_{2,2}=1$). These pieces have dimension $5$,
while by Lemma~\ref{lem:support} and the computation of Section~\ref{sec:checks} the matrices with $2\in K(P)$ lie in
finitely many polytopes of dimension at most $4$. So the condition holds on a set of positive five-dimensional measure
of matrices with $2\notin K(P)$.
\end{remark}

\begin{example}
In Example~3.10 of~\cite{GH26}, $S=\{1,2,3,4\}$ and the support of $\mu$ consists of the maps written $(1212)$, $(1221)$,
$(3434)$ and $(3443)$, where $(j_1j_2j_3j_4)$ denotes the map $r\mapsto j_r$. Here $k(\mu)=2$. The map $e=(1212)$ is an
idempotent of rank two, so $W=\{1,2\}$, $Z_1=\{1,3\}$ and $Z_2=\{2,4\}$. Every map in the support sends $1$ into
$\{1,3\}$ and $2$ into $\{2,4\}$, so $B$ is the identity matrix. In Example~6.2 of~\cite{GH26} the uniform matrix on six
states has a grand coupling with $k=2$; there $|W|=2$ and each $Z_w$ has three elements, and~\eqref{eq:columns} reads
$2\cdot3/6=1$.
\end{example}

\section{A computer check}\label{sec:checks}

The proofs above do not depend on computation. As a check on the statements, a short program~\cite{Repo} enumerates
finite sets $\Sigma$ of maps, computes the semigroup $M(\Sigma)$ and its least rank $r(\Sigma)$, and, whenever $r(\Sigma)\ge2$,
verifies~\eqref{eq:pointwise} for an idempotent of least rank. For $n=3$ it treats every $\Sigma$ with at most three
elements (3303 sets, 233 of them with $r(\Sigma)\ge2$), and for $n=4$ and $n=5$ it treats $10^5$ random sets of one to
four maps each ($25029$ and $24989$ of them with $r(\Sigma)\ge2$). A second program enumerates, for $n\le5$, every
non-empty family $\mathcal H$ of $k$-element subsets of $S$ with $k\ge2$, computes $\operatorname{End}(\mathcal H)$, and
finds the affine dimension of $\{M_f:f\in\operatorname{End}(\mathcal H)\}$ as the rank of an integer matrix. There are
$8$, $79$ and $2078$ such families for $n=3,4,5$, and the largest dimension is $4$, $9$ and $16$ respectively, below
the dimension $n(n-1)$ of the set of stochastic matrices; by Lemma~\ref{lem:support} this gives a second proof of
Theorem~\ref{thm:main} for $n\le5$.

\section*{Acknowledgements and funding}
This work was done by multiple AI agents, Claude (Anthropic) and Codex (OpenAI), connected to
Math Brain~\cite{MathBrainRepo} and working together inside the Brain framework~\cite{BrainFramework}, a shared,
persistent memory in which the agents record claims,
proofs, failed attempts and methods, retrieve one another's results, coordinate who works on what, and audit each
other's work. For more on the Brain and Math Brain, see the Math Brain paper~\cite{MathBrainPaper}, the Math Brain
repository~\cite{MathBrainRepo} and the Brain framework repository~\cite{BrainFramework}.

The research presented in this paper was supported by the European Research Council (ERC) under the European Union's
Horizon Europe research and innovation programme (grant agreement No.~101041711), by the Simons Foundation as part of
the Collaboration on the Mathematical and Scientific Foundations of Deep Learning, by Heights Labs, by Convex Nexus
Capital, by the Israel Science Foundation (grant number 2258/19), and by the Israel Science Foundation (ISF Grant
4101/25). Computational resources were provided in part by Google Cloud through its education credits programme
(credit 530550103).

\section*{Declaration of competing interest}
The author holds equity in Heights Labs and in Convex Nexus Capital and has a management role at Convex Nexus Capital.
These entities are included in the funding declaration above.

\section*{Data availability}
The programs described in Section~\ref{sec:checks} are in the repository~\cite{Repo}.

\section*{Declaration of generative AI and AI-assisted technologies}
This paper is part of a broader research project on AI proving problems. Claude agents (Anthropic) proved
Theorems~\ref{thm:main} and~\ref{thm:structure}, wrote the programs of Section~\ref{sec:checks} and prepared the
manuscript; Codex agents (OpenAI) refereed the proofs. The grand couplings, the coalescence number, the set $K(P)$, the
law $Q$ and the question answered here are due to Grimmett and Holmes~\cite{GH21,GH26}, and Lemma~\ref{lem:rank} is
their Theorem~3(a) of~\cite{GH21}. The permutation property in Lemma~\ref{lem:idempotent} is due to Yano~\cite{Yan13}.
The new results are the linear condition on $P$ in Theorem~\ref{thm:structure} and Theorem~\ref{thm:main}. The author
directed the work
and is responsible for the paper.

\begin{thebibliography}{99}
\bibitem{Bir46} G.~Birkhoff, \emph{Three observations on linear algebra}, Univ. Nac. Tucum\'an. Revista A \textbf{5}
(1946), 147--151.
\bibitem{BM00} G.~Budzban and A.~Mukherjea, \emph{A semigroup approach to the road coloring problem}, in: Probability
on Algebraic Structures, Contemp. Math.~261, Amer. Math. Soc., Providence, RI, 2000, pp.~195--207,
\doi{10.1090/conm/261/04141}.
\bibitem{Fri90} J.~Friedman, \emph{On the road coloring problem}, Proc. Amer. Math. Soc. \textbf{110} (1990),
1133--1135, \doi{10.2307/2047767}.
\bibitem{GH21} G.~R.~Grimmett and M.~Holmes, \emph{Non-coupling from the past}, in: In and Out of Equilibrium~3:
Celebrating Vladas Sidoravicius, Progress in Probability~77, Birkh\"auser/Springer, Cham, 2021, pp.~471--485,
\doi{10.1007/978-3-030-60754-8_22}.
\bibitem{GH26} G.~R.~Grimmett and M.~Holmes, \emph{Coalescence in Markov chains}, J. Theoret. Probab. \textbf{39} (2026), no.~3,
\doi{10.1007/s10959-026-01500-w}, \arxiv{2510.13572}.
\bibitem{How95} J.~M.~Howie, \emph{Fundamentals of Semigroup Theory}, London Mathematical Society Monographs, New
Series~12, Clarendon Press, Oxford, 1995.
\bibitem{PW96} J.~G.~Propp and D.~B.~Wilson, \emph{Exact sampling with coupled Markov chains and applications to
statistical mechanics}, Random Structures Algorithms \textbf{9} (1996), 223--252.
\bibitem{Yan13} K.~Yano, \emph{Random walk in a finite directed graph subject to a road coloring}, J. Theoret.
Probab. \textbf{26} (2013), 259--283, \doi{10.1007/s10959-011-0398-8}, \arxiv{1005.0079}.
\bibitem{Repo} Y.~Shmalo, \emph{Coalescence of grand couplings: paper and check programs}, 2026,
\url{https://github.com/yspennstate/grand-coupling-coalescence}.
\bibitem{MathBrainRepo} Y.~Shmalo, \emph{Math Brain: public snapshot, code and tests}, 2026,
\url{https://github.com/yspennstate/math-brain-public/tree/8c9a6bd478c4ec4082347151b7ae16db50c917dd}.
\bibitem{MathBrainPaper} Y.~Shmalo, \emph{Math Brain: shared, auditable memory for multi-agent mathematical research},
2026, \doi{10.5281/zenodo.24150837},
\url{https://github.com/yspennstate/math-brain-public/tree/main/papers/math-brain-2026-10-24}.
\bibitem{BrainFramework} Y.~Shmalo, \emph{Brain Framework}, public software repository, 2026,
\url{https://github.com/yspennstate/brain-framework-public/tree/5629e644cdd510173ae97e9aa92642dabd22dc57}.
\end{thebibliography}

\end{document}
